% Generated draft: requires mathematical and source review.
\documentclass[11pt]{article}
\usepackage[margin=25mm]{geometry}
\usepackage{fontspec}
\setmainfont{Latin Modern Roman}
\usepackage{amsmath,amssymb,mathtools,unicode-math}
\setmathfont{Latin Modern Math}
\usepackage{enumitem,needspace,xurl,hyperref}
\setlength{\parindent}{0pt}
\setlength{\parskip}{5pt}
\setlength{\emergencystretch}{3em}
\newcommand{\N}{\mathbb{N}}
\newcommand{\Z}{\mathbb{Z}}
\newcommand{\Q}{\mathbb{Q}}
\newcommand{\R}{\mathbb{R}}
\newcommand{\C}{\mathbb{C}}
\newcommand{\F}{\mathbb{F}}
\newcommand{\A}{\mathbb{A}}
\newcommand{\PP}{\mathbb{P}}
\newcommand{\E}{\mathbb{E}}
\newcommand{\eps}{\varepsilon}
\newcommand{\dd}{\,\mathrm{d}}
\begin{document}
\section*{UnsolvedMath ID: 20000232}
\textbf{Source problem code:} catalogue entry
\subsection{Definitions and mathematical statement}
The AIM question asks for a new proof of Green's theorem by tropical methods. The attached report interprets this as the generic canonical-syzygy theorem: for a general smooth projective complex curve \(C\) of genus \(g\), the canonical Koszul groups vanish in the Green range \(K_{p,2}(C,K_C)=0\) for \(p<\operatorname{Cliff}(C)\), with generic Clifford index \(\lfloor(g-1)/2\rfloor\). Here \(\operatorname{Cliff}(C)\) is the minimum of \(\deg L-2(h^0(L)-1)\) over special line bundles with \(h^0(L),h^1(L)\geq2\), when applicable. For \(V=H^0(C,K_C)\), \(K_{p,2}\) is the middle cohomology of the canonical Koszul complex
\[
\bigwedge^{p+1}V\otimes H^0(K_C)
\longrightarrow\bigwedge^pV\otimes H^0(K_C^2)
\longrightarrow\bigwedge^{p-1}V\otimes H^0(K_C^3).
\]
The exact formulation and low-genus conventions intended by ``Green's theorem'' are not specified in the original. The report says the generic theorem is known algebraically but no all-genus tropical proof was located; the latter is a report of its search, not a new literature determination.

Its candidate obstruction asserts that for every integer \(g\geq5\), there are two constant good-reduction curves over the same valued field \(\C((t))\) with the same one-vertex weighted skeleton and tropical canonical divisor, but with
\[
K_{1,2}(C_0,K_{C_0})=0,\qquad K_{1,2}(C_1,K_{C_1})\ne0.
\]
For that skeleton the unique vertex \(v\) has weight \(g\) and its canonical divisor is \((2g-2)v\). Thus abstract weighted-skeleton divisor data alone do not determine even this first nonlinear canonical syzygy group. This does not rule out tropical proofs using enriched degeneration or initial-ideal information.

The report separately claims an explicit flat three-quadric family in \(\PP^4\) whose \(t\)-adic initial canonical ideal is
\[
\operatorname{in}_t(I)=(x_0^2,x_1^2,x_2^2),
\]
whose generic fiber is a smooth canonical genus-five curve, and whose complete-intersection Koszul resolution certifies the critical vanishing \(K_{1,2}=0\). In the stated complete-intersection situation the resolution has the shape
\[
0\longrightarrow S(-6)\longrightarrow S(-4)^3
\longrightarrow S(-2)^3\longrightarrow S\longrightarrow S/I\longrightarrow0,
\]
where \(S\) is the five-variable graded polynomial ring over the generic field. The supplied summary does not contain the claimed explicit three quadrics or their flatness and smoothness certificate. Consequently this records the proposed special-case certificate and its initial ideal, not a reproducible explicit family or an all-genus tropical proof.
\subsection{Short English statement}
Seek a tropical proof of the generic Green syzygy theorem. The report proposes that abstract skeleton data are insufficient, and offers a genus-five initial-ideal example using richer degeneration data, without claiming an all-genus proof.
\subsection{Sources}
\begin{itemize}
\item AIM workshop ``Degenerations in algebraic geometry,'' Problem 8.3: \url{https://aimath.org/pastworkshops/degenalggeomproblems.pdf}. The PDF was not accessed here.
\item Local canonical registry \texttt{aim-algebraic-geometry-notes.json}, \texttt{notes[231]}, tag \texttt{problem}; permanent UnsolvedMath ID \(20000232\), source code \texttt{catalogue entry}. Exact original asks: ``Can one give a new proof of Green's theorem by tropical methods?'' The trailing extracted ``49'' has no specified mathematical role and is treated as a layout artifact. Original remarks and literature field are empty.
\item Supplied research attempt 1 proposing the skeleton obstruction and genus-five initial-ideal certificate. No full authoritative bibliographic references for the algebraic theorem, tropical status, or candidate examples are supplied. The report is machine generated, not peer reviewed, and requires independent verification before citation.
\end{itemize}
\end{document}
